\section{Derive a shuffle across an input boundary}
Number the 48 input bytes of a sixteen-pixel group from zero through 47. Pixel four
uses red byte 12, green 13 and blue 14. Pixel five uses red 15, green 16 and blue 17,
so it straddles the boundary between the first and second input vectors. The second
output vector must contain words for pixels four through seven.

\Needspace{9\baselineskip}
\begin{longtable}{p{24mm}p{42mm}p{35mm}p{44mm}}
  \toprule Pixel & Global RGB byte indices & Output byte-lane sources & Shuffle selectors \\\midrule
  4              & 12, 13, 14              & 14, 13, 12, zero         & 14, 13, 12, 128   \\
  5              & 15, 16, 17              & 17, 16, 15, zero         & 17, 16, 15, 128   \\
  6              & 18, 19, 20              & 20, 19, 18, zero         & 20, 19, 18, 128   \\
  7              & 21, 22, 23              & 23, 22, 21, zero         & 23, 22, 21, 128   \\\bottomrule
\end{longtable}
For this shuffle, put the first input vector on the right and the second on the
left. Selectors zero through fifteen select bytes of the right operand; sixteen
through thirty-one select bytes of the left operand. Selector 17 therefore chooses
lane one of the second vector, which is global byte 17. A selector with bit seven
set produces zero. Store its byte pattern as \texttt{-128} in a signed
\texttt{.byte} declaration; the control bits are hexadecimal \texttt{80}.

After the complete-group bounds checks, let \texttt{\$a0} point to sample byte zero,
\texttt{\$t0} to output word zero, and \texttt{\$w5} hold this second selector vector.
The following excerpt writes output words four through seven:
\MARSCodeLabel{One output vector spanning two input vectors}
\begin{minted}[samepage=true]{mars-idea}
ld.b $w0, 0($a0)
ld.b $w1, 16($a0)
move.v $w3, $w5
vshf.b $w3, $w1, $w0
st.w $w3, 16($t0)
\end{minted}
The destination initially supplies the selectors to \texttt{vshf.b}; afterward it
contains shuffled data. Copying the retained mask into \texttt{\$w3} makes the mask
available again on the next iteration. The four byte lanes B, G, R, zero form the
numeric word \texttt{0x00RRGGBB}. The word store writes that value using the target's
memory byte order. For \texttt{0x00112233}, the memory bytes are
\texttt{33 22 11 00} in little-endian order and \texttt{00 11 22 33} in big-endian
order. The displayed word value is the same in both cases.

Derive the other three masks from the same global indices. The third output vector
spans input vectors one and two, so subtract sixteen from its global byte indices
before choosing between those operands. The fourth output vector uses only the last
input vector. A complete group consumes 48 sample bytes and publishes 64 pixel bytes;
advance each cursor by its own amount before testing for another sixteen pixels.
